\section{Encoder 预训练：BERT 与掩码语言建模}

\subsection{三种 Transformer 架构回顾}

预训练方法与 Transformer 的架构类型密切相关。Chris Manning 将模型分为三类：

\begin{figure}[H]
\centering
\includegraphics[width=0.95\textwidth]{slides-images/slide-030.jpg}
\caption{三种预训练架构：Encoder（双向上下文，用于表示学习）、Encoder-Decoder（编码+生成）、Decoder（单向上下文，用于语言建模/生成）\protect\footnotemark}
\end{figure}
\footnotetext{Slides 第31页。}

\begin{knowledgebox}{三种架构的特点与适用场景}
\begin{itemize}
    \item \textbf{Encoder}：每个 token 能看到整个序列（双向上下文）。适合\textbf{理解}类任务（分类、NER、问答抽取）。代表模型：BERT
    \item \textbf{Encoder-Decoder}：Encoder 部分看到完整输入，Decoder 部分自回归生成输出。适合\textbf{序列到序列}任务（翻译、摘要）。代表模型：T5
    \item \textbf{Decoder}：每个 token 只能看到它之前的 token（单向上下文）。适合\textbf{生成}类任务。代表模型：GPT 系列
\end{itemize}
\end{knowledgebox}

\subsection{为什么 Encoder 不能用语言建模预训练？}

Encoder 拥有双向上下文，这意味着在预测下一个词时，模型已经\textbf{看到了}那个词——预测任务变得毫无意义（trivial）。

\begin{figure}[H]
\centering
\includegraphics[width=0.85\textwidth]{slides-images/slide-031.jpg}
\caption{Encoder 中的语言建模问题：由于双向注意力，位置 $i$ 的表示已经包含了位置 $i+1$ 的信息，``预测下一个词''变成了简单的复制操作\protect\footnotemark}
\end{figure}
\footnotetext{Slides 第32页。}

解决方案：\textbf{掩码语言建模}（Masked Language Modeling, MLM）——先遮住一些词，再让模型预测被遮住的词。

\subsection{BERT：掩码语言建模}

BERT（Bidirectional Encoder Representations from Transformers，2018）是第一个成功将掩码语言建模大规模应用于预训练的模型。

\begin{figure}[H]
\centering
\includegraphics[width=0.95\textwidth]{slides-images/slide-032.jpg}
\caption{掩码语言建模示意：输入 ``I [MASK] to the [MASK]''，模型利用双向上下文预测 ``went'' 和 ``store''。损失仅在被遮住的位置计算\protect\footnotemark}
\end{figure}
\footnotetext{Slides 第33页。}

BERT 的掩码策略选取 15\% 的 token 进行处理，对于被选中的 token：

\begin{itemize}
    \item \textbf{80\%} 替换为 \texttt{[MASK]} 标记
    \item \textbf{10\%} 替换为词表中的随机词
    \item \textbf{10\%} 保持不变
\end{itemize}

\begin{warningbox}{为什么不能全部用 [MASK]？}
如果只在 \texttt{[MASK]} 标记处做预测，模型可能会学到：``只有遇到 [MASK] 时才需要认真构建表示，其他位置可以偷懒''。但在实际使用时（如情感分析），输入中\textbf{不会出现} [MASK] 标记。通过混合使用随机替换和保持不变，模型被迫在\textbf{任何位置}都构建高质量的上下文表示。
\end{warningbox}

\begin{figure}[H]
\centering
\includegraphics[width=0.95\textwidth]{slides-images/slide-034.jpg}
\caption{BERT 的输入表示：Token Embedding + Segment Embedding + Position Embedding，三者\textbf{逐元素相加}（非拼接）\protect\footnotemark}
\end{figure}
\footnotetext{Slides 第35页。来自 Devlin et al., 2018。}

\subsection{下一句预测任务（NSP）}

BERT 原始论文还引入了一个辅助任务——\textbf{下一句预测}（Next Sentence Prediction, NSP）：给定两段文本 A 和 B，预测 B 是否是 A 的真实续文。

这个任务的设计初衷是教会模型理解\textbf{长距离的文本连贯性}。然而，后续研究（RoBERTa，2019）发现这个任务\textbf{并不必要}，原因有二：

\begin{enumerate}
    \item NSP 将有效上下文长度减半（原本 512 token 的窗口被分成两段各 256 token）
    \item 模型在 NSP 任务上表现很差，可能是因为这个任务对当时的模型来说太难了
\end{enumerate}

\subsection{BERT 在 GLUE 基准上的表现}

BERT 的发布在 NLP 领域引起了巨大震动。在 GLUE 基准（包含多种 NLP 任务）上，BERT 全面超越了之前\textbf{为每个任务精心设计}的专用模型。

\begin{figure}[H]
\centering
\includegraphics[width=0.95\textwidth]{slides-images/slide-035.jpg}
\caption{BERT 在 GLUE 基准上的表现：BERT$_{\text{LARGE}}$ 在所有 8 个任务上均超越了此前的最好成绩，平均分从 74.0 提升到 82.1\protect\footnotemark}
\end{figure}
\footnotetext{Slides 第36页。来自 Devlin et al., 2018。}

Chris Manning 评价道：``All of that was blown out of the water by just build a big Transformer and teach it to predict the missing words.'' BERT 的成功宣告了``为每个任务设计专用架构''时代的终结。

\begin{importantbox}{BERT 的核心规格}
\begin{itemize}
    \item \textbf{BERT$_{\text{BASE}}$}：110M 参数，12 层，768 维隐层，12 头
    \item \textbf{BERT$_{\text{LARGE}}$}：340M 参数，24 层，1024 维隐层，16 头
    \item \textbf{训练数据}：BookCorpus + English Wikipedia（约 25 亿词）
    \item \textbf{预训练计算}：在当时被认为``只有 Google 才负担得起''，但今天看来并不算大
    \item \textbf{微调}：在单个 GPU 上即可完成
\end{itemize}
\end{importantbox}

\subsection{BERT 的微调方式}

将预训练好的 BERT 用于下游任务非常简单：

\begin{figure}[H]
\centering
\includegraphics[width=0.85\textwidth]{slides-images/slide-036.jpg}
\caption{BERT 微调示意：移除预训练的预测头，在 BERT 输出之上添加一个新的线性分类器，然后端到端微调整个网络\protect\footnotemark}
\end{figure}
\footnotetext{Slides 第37页。}

\begin{enumerate}
    \item 移除预训练时的 MLM 预测头
    \item 对于分类任务：取 \texttt{[CLS]} 标记或最后一个 token 的输出向量，接一个线性分类器
    \item 对于 token 级别任务（如 NER）：对每个 token 的输出向量接线性分类器
    \item 整个网络（BERT + 新分类器）一起微调
\end{enumerate}

\subsection{BERT 的局限性}

\begin{warningbox}{BERT 不适合文本生成}
BERT 的双向上下文使其非常擅长``理解''任务（分类、抽取式问答、NER），但它\textbf{没有自回归生成的能力}——无法像语言模型那样逐词生成文本。如果你需要生成摘要、翻译或对话回复，应该使用 Encoder-Decoder 或 Decoder-only 模型。
\end{warningbox}

\subsection{BERT 的改进：SpanBERT 与 RoBERTa}

\textbf{SpanBERT}：用\textbf{连续片段掩码}（span masking）替代随机单词掩码。随机掩码单个子词过于简单（前后子词的线索太强），而掩码一整个连续片段迫使模型进行更深层的理解。

\textbf{RoBERTa}（Robustly Optimized BERT Approach）：
\begin{itemize}
    \item 移除了 NSP 任务
    \item 在更多数据上训练更长时间
    \item 动态掩码（每次 epoch 重新生成掩码位置）
    \item 是 BERT 的\textbf{直接替代品}，性能更优
\end{itemize}

\begin{knowledgebox}{选择建议}
如果你的项目需要使用 encoder 模型（如分类、NER、抽取式问答），优先使用 \textbf{RoBERTa} 而非原始 BERT，它在几乎所有任务上都表现更好，且完全兼容 BERT 的使用方式。
\end{knowledgebox}

\subsection{参数高效微调}

当模型变得越来越大时，全参数微调变得昂贵且不实际。\textbf{参数高效微调}（Parameter-Efficient Fine-Tuning, PEFT）只调整少量参数，而冻结大部分预训练参数。

\begin{figure}[H]
\centering
\includegraphics[width=0.85\textwidth]{slides-images/slide-038.jpg}
\caption{Prefix Tuning 示意：冻结所有预训练参数（蓝色），仅训练前缀向量（红色）——一组``伪词向量''被拼接到输入序列最前面\protect\footnotemark}
\end{figure}
\footnotetext{Slides 第39页。}

三种主要的参数高效微调方法：

\begin{enumerate}
    \item \textbf{Prefix Tuning / Prompt Tuning}：冻结网络，在序列前添加可训练的``伪词向量''
    \item \textbf{Adapter Layers}：在 Transformer 层之间插入小型可训练的适配层
    \item \textbf{LoRA}（Low-Rank Adaptation）：冻结原始权重矩阵 $W$，学习一个低秩增量 $\Delta W = BA$，令 $W' = W + BA$
\end{enumerate}

\begin{importantbox}{LoRA 的核心思想}
对于权重矩阵 $W \in \mathbb{R}^{d \times d}$，LoRA 学习两个小矩阵 $B \in \mathbb{R}^{d \times r}$ 和 $A \in \mathbb{R}^{r \times d}$（$r \ll d$），令实际权重为：
$$W' = W + BA$$
只训练 $B$ 和 $A$（参数量从 $d^2$ 降至 $2dr$），冻结 $W$。这背后的直觉是：从预训练到下游任务的参数变化是\textbf{低秩的}，不需要修改全部参数。
\end{importantbox}

\subsection{本章小结}

Encoder 预训练以 BERT 为代表，通过掩码语言建模让模型学会双向上下文表示。BERT 在各类 NLP 基准上的巨大成功标志着``预训练--微调''范式的确立。后续的 SpanBERT、RoBERTa 在训练策略上进行了优化。参数高效微调方法（Prefix Tuning、LoRA）则解决了大模型微调的效率问题。但 Encoder 模型不适合生成任务，这引出了 Encoder-Decoder 和 Decoder-only 架构。

%% ========== Part 4: Encoder-Decoder 预训练 ==========

